% ============================================================
% 30_layout_spacing.tex — Layout-Grundlagen, Spacing-Skala, Box-Bindungen
% ============================================================

% MaschKo-Verdichtung: kompakteres parskip (6-Seiten-Budget, bildlastige ZSF)
\setlength{\parskip}{0.5pt plus 0.8pt minus 0.2pt}
% MaschKo-Verdichtung: engere Zeilenabstände (Inspiration nutzte 6.75pt/6pt;
% 0.86 für das 6-Seiten-Ziel — immer noch über der Vorlage (6.75/6 = 0.83)
\linespread{0.74} % 6-Seiten-Ziel: baselineskip ~6.0pt wie die Inspirationsfassung
\setlength{\parindent}{0pt}

\widowpenalty=10000
\clubpenalty=10000
\brokenpenalty=4000
\predisplaypenalty=10000

\raggedcolumns
\newcounter{zsfFirstChap}\setcounter{zsfFirstChap}{1}

% ------------------------------------------------------------
% SPACING-SKALA (4 Stufen) — die einzige Quelle für Abstände
% ------------------------------------------------------------
% MaschKo-Verdichtung: Skala enger als Template-Default (4/7/10)
\newlength{\ZSFspaceXXS}\setlength{\ZSFspaceXXS}{0.5pt} % MaschKo: Mikro-Stufe für Display-Skips (6-Seiten-Ziel)
\newlength{\ZSFspaceXS}\setlength{\ZSFspaceXS}{1pt}
\newlength{\ZSFspaceS} \setlength{\ZSFspaceS}{2pt}
\newlength{\ZSFspaceM} \setlength{\ZSFspaceM}{4pt}
\newlength{\ZSFspaceL} \setlength{\ZSFspaceL}{6pt}

% Outer (Text-Box-Bindung) und Sep (Trennlinien)
\newlength{\ZSFspaceOuter}\setlength{\ZSFspaceOuter}{2pt} % MaschKo-Verdichtung: 4 -> 2 (6-Seiten-Ziel)
\newlength{\ZSFspaceSep}  \setlength{\ZSFspaceSep}{1pt} % MaschKo-Verdichtung: 3 -> 1.5 -> 1 (Trennstriche frassen ~8pt pro \formulasep)

% Corner-Radius-Tokens: runder Grundlook, aber zentral steuerbar.
\newlength{\ZSFcornerBox}   \setlength{\ZSFcornerBox}{0.6mm}
\newlength{\ZSFcornerInline}\setlength{\ZSFcornerInline}{1mm}

% ------------------------------------------------------------
% Spalten-Layout
% ------------------------------------------------------------
\setlength{\columnsep}{4pt} % MaschKo-Verdichtung: L (6pt) -> 4pt (6-Seiten-Ziel)
\setlength{\multicolsep}{\ZSFspaceS}
\setlength{\emergencystretch}{2em}

% ------------------------------------------------------------
% Schwellwerte und Titelbalken-Hilfsgrössen
% ------------------------------------------------------------
\newlength{\ZSFbreakThreshold}\setlength{\ZSFbreakThreshold}{0pt} % Schwellwert minimiert/deaktiviert: verhindert vorzeitige Spaltenumbrüche bei geringem Restplatz
\newlength{\ZSFbreakThresholdChapter}\setlength{\ZSFbreakThresholdChapter}{3\baselineskip} % Mindesthöhe für ChapterBar (verhindert Bars am Spaltenende); MaschKo: 3 statt 10 (Platzbudget; Boxen sind breakable, Bar+erste Teilbox brauchen weniger Reserve)
\newlength{\ZSFpostChapterNudge}\setlength{\ZSFpostChapterNudge}{0pt}
\newlength{\ZSFtitleBarHeight}\setlength{\ZSFtitleBarHeight}{5.8pt} % MaschKo: 9.2 -> 7 -> 5.8 (6-Seiten-Ziel; Strut der Titelschrift bestimmt die effektive Höhe)
\newlength{\ZSFsubsectionBarBeforeSkip}\setlength{\ZSFsubsectionBarBeforeSkip}{4pt} % MaschKo: mehr Luft vor echten Unterthemen, bleibt im 6-Seiten-Budget
\newlength{\ZSFsubsubsectionBarBeforeSkip}\setlength{\ZSFsubsubsectionBarBeforeSkip}{\ZSFspaceXXS} % dritte Ebene: so kompakt wie bisherige Stern-Zwischentitel
\newlength{\ZSFsubsectionAfterChapterGap}\setlength{\ZSFsubsectionAfterChapterGap}{0pt} % Kapitelbalken + erster Abschnitt ohne Zwischenlücke
\newlength{\ZSFsubsectionAfterGap}\setlength{\ZSFsubsectionAfterGap}{\ZSFspaceXS}
\newlength{\ZSFchapterBarAfterSkip}\setlength{\ZSFchapterBarAfterSkip}{1pt} % MaschKo-Verdichtung: 2 -> 1 (6-Seiten-Ziel)

% Maximale Höhe für Bilder in figboxen (verhindert überhohe Aspect-Ratios).
\newlength{\ZSFfigMaxHeight}\setlength{\ZSFfigMaxHeight}{2.6cm}
% Globaler Skalierungsfaktor für \ZSFfig-Bilder (Seitenbudget-Hebel).
% Getestet (2026-07): 0.93 reduziert die Seitenzahl NICHT (Packverlust
% dominiert) — deshalb 1.0 = exakte Vorlagen-Verhältnisse. Werte bis 0.90
% liegen innerhalb der 10%-Toleranz, falls später nötig.
\newcommand{\ZSFfigScale}{1.0}

% Tokens für statementbox / procedure
\newlength{\ZSFstatementBarWidth}\setlength{\ZSFstatementBarWidth}{2pt}
\newlength{\ZSFstatementLeftPad} \setlength{\ZSFstatementLeftPad}{\ZSFspaceM}
\newlength{\ZSFprocStepGap}      \setlength{\ZSFprocStepGap}{\ZSFspaceXS}
\newlength{\ZSFprocItemSep}      \setlength{\ZSFprocItemSep}{\ZSFspaceXS}

% Padding für cellspace (siehe styles/20_tables.tex Konvention)
\newlength{\ZSFtextMinPad}\setlength{\ZSFtextMinPad}{\ZSFspaceXS}

% ------------------------------------------------------------
% Code-Längen — nur für die optionalen Code-Module benötigt
% (styles/65_code_style.tex + 67_code_comments.tex). Schadlos
% wenn die Module in preamble.tex deaktiviert sind.
% MaxWidth:   max. Breite des rechts-stehenden Kommentars
% Threshold:  Schwelle für Auto-Overflow (Code+Kommentar > Threshold -> oben)
% Default-Register sind immutabel und dienen \ResetCodeCommentThreshold
% als Anker; \SetCodeCommentThreshold modifiziert nur die "lebenden" Register.
% ------------------------------------------------------------
\newlength{\ZSFcodeMiniSkip}  \setlength{\ZSFcodeMiniSkip}{1mm}
\newlength{\ZSFcodeFrameRule} \setlength{\ZSFcodeFrameRule}{0.4pt}
\newlength{\ZSFcodeXMargin}   \setlength{\ZSFcodeXMargin}{8pt}
\newlength{\ZSFcodeCommentMaxWidthDefault}
  \setlength{\ZSFcodeCommentMaxWidthDefault}{.25\columnwidth}
\newlength{\ZSFcodeCommentThresholdDefault}
  \setlength{\ZSFcodeCommentThresholdDefault}{\maxdimen}
\newlength{\ZSFcodeCommentMaxWidth} \setlength{\ZSFcodeCommentMaxWidth}{\ZSFcodeCommentMaxWidthDefault}
\newlength{\ZSFcodeCommentThreshold}\setlength{\ZSFcodeCommentThreshold}{\ZSFcodeCommentThresholdDefault}
\newlength{\ZSFcodeCommentGap}      \setlength{\ZSFcodeCommentGap}{\ZSFspaceM}

% ============================================================
% TEXT-BOX-BINDUNGEN
% Koppeln Fliesstext direkt an Boxen — verhindert visuellen Bruch.
% ============================================================

% Robustes Skip-/Penalty-Aufrollen (vmode-/hmode-aware)
\newcommand{\ZSFRobustUnskip}{%
  \ifvmode
    \ifdim\lastskip>0pt\relax\vskip-\lastskip\fi
    \ifnum\lastpenalty=0\else\unpenalty\fi
    \ifdim\lastskip>0pt\relax\vskip-\lastskip\fi
  \else
    \ifdim\lastskip>0pt\relax\unskip\fi
    \ifnum\lastpenalty=0\else\unpenalty\fi
    \ifdim\lastskip>0pt\relax\unskip\fi
  \fi
}

% Text der zur Box DARUNTER gehört (Einleitung)
\newcommand{\textVorBox}[1]{%
  \par
  \addvspace{\ZSFspaceOuter}%
  \begingroup
  \setlength{\parskip}{0pt}%
  {\small\color{TextVorBoxColor} \noindent #1\par}%
  \endgroup
  \nopagebreak
  \vspace{\ZSFspaceXS}%
  \vspace{-\ZSFspaceS}%
  \nopagebreak
}

% Text der zur Box DARÜBER gehört (Ergänzung/Folgerung)
\newcommand{\textNachBox}[1]{%
  \par
  \ZSFRobustUnskip
  \vspace{\ZSFspaceXS}%
  \begingroup
  \setlength{\parskip}{0pt}%
  \nopagebreak
  {\small\color{TextNachBoxColor} \noindent #1\par}%
  \endgroup
  \addvspace{\ZSFspaceOuter}%
}

% Titel/Text vor einer Formel (innerhalb der formulabox)
\newcommand{\textVorFormel}[1]{%
  \par
  \begingroup
  \setlength{\parskip}{0pt}%
  \noindent\raggedright\textcolor{black!60}{\textbf{#1}}\par
  \endgroup
  \vspace{\ZSFspaceXS}%
  \centering
}

% Erklaerungstext nach einer Formel (innerhalb der formulabox)
\newcommand{\textNachFormel}[1]{%
  \par
  \vspace{\ZSFspaceXS}%
  \begingroup
  \setlength{\parskip}{0pt}%
  \noindent\raggedright\textcolor{black!60}{\small #1}\par
  \endgroup
  \centering
}

% ------------------------------------------------------------
% Semantische Gap-Helfer für Kapitel (sparsam einsetzen)
% ------------------------------------------------------------
\newcommand{\ZSFSectionGap}{\vspace{\ZSFspaceL}}
\newcommand{\ZSFgapXS}{\vspace{\ZSFspaceXS}}
\newcommand{\ZSFgapS}{\vspace{\ZSFspaceS}}
\newcommand{\ZSFgapM}{\vspace{\ZSFspaceM}}
\newcommand{\ZSFgapL}{\vspace{\ZSFspaceL}}
\newcommand{\ZSFmanualColumnBreak}{\columnbreak}
\newcommand{\ZSFtabColGap}{\hspace{1.5em}}
